\chapter{Conclusion and Future Work}
\label{ch:conclusion}

\section{Conclusion}
\label{sec:conclusion}

This thesis asked whether a decision maker for an autonomous factory robot is better built from
explicit rules, from learning, or from a combination in which each part does what it is good at.

The system built to answer that question is complete rather than illustrative: one description of
the hall generates the simulated world, the navigation map and the distances the models reason
with; a factory of three unequal production lines loses production when a buffer fills; a robot
with a modelled battery, motor temperature and drivetrain wear executes whole actions and recovers
when navigation fails; and a decision service, which the robot survives without, answers the
question \emph{what next?} before every action. The same policy object runs in a fast simulator and
on the robot, which is what makes thousands of measured shifts comparable to the handful driven in
full physics.

Four answers came out of it. First, the \NS is the strongest policy that never breaks a constraint:
\num{53.81} against \num{52.91} for its own symbolic layer, \num{52.35} for a deep reinforcement
learning policy with the same information, and \num{41.66} for a rule-based reference, over
\num{3240} shifts. Second, the two halves are not interchangeable. The rules carry most of the
competence and all of the guarantee; the network adds a small, consistent improvement --- and
without the rules, the same network delivers more units than the full model while entering an
unsafe state in one shift in six. It learned the judgement and not the guarantee. Third, that
learned judgement is only useful while it is contained: unbounded, it dragged the model below its
own rule layer on factories it had never seen, and bounding its influence to a single score unit
--- a change selected on training data alone --- recovered \num{3.87} points on unseen scenarios and
\num{21.11} on the hardest of them. Fourth, none of this is an artefact of the fast simulator:
energy transfers to the full Gazebo stack within \SI{0.2}{\percent}, and in a sixty-minute shift
the reinforcement learning policy drove the battery to \SI{9.4}{\percent} while the \NS never went
below \SI{17.0}{\percent}.

The single sentence worth taking away is this: \emph{a learned component improves a rule-based
policy only as long as it cannot overrule it} --- and constraining how much influence it may
exert, rather than only which actions it may choose, is what makes that improvement survive
outside the situations it was trained on.

Two of the three findings in this work came from a measurement that contradicted a decision already
taken: the learned ranker that ignored its own rules, and the path planner chosen on a benchmark
that measured the wrong regime. Both were found by watching complete shifts rather than by
inspecting aggregate scores, which is perhaps the most transferable lesson here.

\section{Future work}
\label{sec:future-work}

The open questions below follow from the limitations of \cref{sec:limitations}; each is stated with
the smallest experiment that would settle it.

\paragraph{Charge the objective for all the energy.}
Idle power drawn while docked is currently paid by the charging station and not counted
(\cref{sec:objective}), which flatters policies that park. Adding it changes what every model is
optimising, so both learned models would have to be retrained and the full evaluation repeated ---
a day of computation, and the most likely single change to alter the ranking of the policies.

\paragraph{Make the bound state-dependent.}
The bound $\beta$ is a single constant, selected once. A natural extension is to widen it where the
network's input is well covered by its training data and to narrow it where it is not, which turns
the fixed trust region into a measure of the model's own competence. The experiment is the one in
\cref{tab:bound-sweep}, repeated with a density or ensemble-disagreement estimate in place of the
constant.

\paragraph{Teach the model to batch.}
The one scenario in which a simpler policy still wins is \texttt{worn\_robot}, where the rule
reference makes fewer, fuller trips (\cref{sec:limitations}). The look-ahead horizon used for
training is \SI{900}{\second}; lengthening it so that the benefit of waiting for a fuller load
appears in the training targets is a cheap experiment with a clear prediction.

\paragraph{Learn the symbolic layer.}
The rules are hand-written, and H5 demonstrated that a necessary one can be missing until a real
shift exposes it. Inducing rules from logged decisions would remove that dependence on the author's
foresight --- provided the induced rules stay readable enough to be checked by a person, which is
the property that makes the current layer valuable.

\paragraph{Beyond one robot, and beyond simulation.}
Several robots sharing one charger and one delivery point turn this into a scheduling problem with
congestion and contention, requiring rules about mutual exclusion and an observation that includes
the other robots: the largest step away from what was built here, and the most industrially
relevant. Everything here is also simulation, and agreement between the two simulators is
consistency rather than proof (\cref{sec:limitations}); a physical AMR would test the energy model
against rolling resistance, wheel slip and a real battery.
